% !PRINTED 275
\weberpage{275}
\weq{5}{P_m, \; P_{m-1}, \; P_{m-2} \ldots P_0}

have, in the interval from $-1$ to $+1$, the properties of a \gesperrt{Sturm} chain. For if $P_n$ and $P_{n-1}$ were to vanish simultaneously for any $x$, then by (2) $P_{n-2}$, and consequently $P_{n-3}$ etc. down to $P_0$, would also have to vanish, which is impossible, since $P_0=1$. Accordingly the condition \S.\,86, 2. is also satisfied for every arbitrary interval.

If $P_{n-1}=0$, then by (2) $P_n$ and $P_{n-2}$ are of opposite sign, as the condition \S.\,86, 3. requires.

From (3) it follows that $P_n(x)$ and $P_n'(x)$ never vanish simultaneously, since otherwise $P_{n-1}(x)$ would also have to vanish; and finally it follows from (3) that, if $P_n(x)$ vanishes and $-1<x<1$, then $P_n'(x)$ and $P_{n-1}(x)$ have the same sign.

If we now substitute the values $-1, +1$ for $x$ in (5), it follows from (4) that for $x=-1$ there are nothing but sign changes in (5), and for $x=+1$ no sign change at all; and from this it follows:

\gesperrt{The equation $P_m(x)=0$ has $m$ real roots between $x=-1$ and $x=+1$, and since $P_m$ is of degree $m$, these are all the roots.}

But we can draw a conclusion that goes somewhat further still. If we take an interval $\varDelta=(\alpha,\beta)$ in which two roots $\xi, \eta$ of $P_m=0$, and no more, are contained, then in the passage from $\alpha$ to $\beta$ two sign changes are lost in the chain (5). We take $\alpha$ so close to $\xi$, and $\beta$ so close to $\eta$, that $P_{m-1}(\alpha)$ has the same sign as $P_{m-1}(\xi)$, and $P_{m-1}(\beta)$ the same sign as $P_{m-1}(\eta)$. If $P_m$ passes from negative to positive as it goes through $\xi$, then it passes from positive to negative as it goes through $\eta$; hence $P_m'(\xi)$, and consequently $P_{m-1}(\xi)$ and $P_{m-1}(\alpha)$, are positive, while $P_m'(\eta)$, and consequently $P_{m-1}(\eta)$ and $P_{m-1}(\beta)$, are negative; and the reverse takes place if $P_m$ passes from positive to negative as it goes through $\xi$. One sees from this, then, that in the passage from $\alpha$ to $\beta$, in

\[
P_m, \quad P_{m-1}
\]

one sign change is lost; hence in the same passage, and consequently also in the passage from $\xi$ to $\eta$, there must, in the chain

\[
P_{m-1}, \quad P_{m-2} \ldots P_0
\]
%CONT>>
